\chapter{Methodology}

\section{System Design}

\subsection{Embedded Subsystem (ESP32)}

\subsubsection{Heart Rate}

Heart rate is derived from the photoplethysmography (PPG) infrared signal by detecting
peaks in the waveform — each peak corresponding to one cardiac cycle.

\subsubsection{SpO\textsubscript{2}}

Blood oxygen saturation is estimated by pulse oximetry, exploiting the difference in
optical absorption between oxyhaemoglobin (HbO\textsubscript{2}) and deoxyhaemoglobin
(Hb) at two wavelengths.

\subsubsection{Delta-Threshold Upload Filter}

Records are enqueued only when the measured value has changed by more than a defined
threshold since the last successful push.
Let $x_n$ be the current reading and $\hat{x}$ the last transmitted value.
The upload condition for room data is:

\begin{equation}
	\text{Upload}_{\text{room}} \;\Leftrightarrow\;
	\left| T_n - \hat{T} \right| > \delta_T \;\lor\;
	\left| H_n - \hat{H} \right| > \delta_H
	\label{eq:room_thresh}
\end{equation}

with $\delta_T = 0.3$\,°C and $\delta_H = 0.5$\,\%.
For patient data:

\begin{equation}
	\text{Upload}_{\text{patient}} \;\Leftrightarrow\;
	\text{BPM}_n \neq \hat{\text{BPM}} \;\lor\;
	\left| \text{SpO}_{2,n} - \widehat{\text{SpO}_2} \right| > \delta_S
	\label{eq:patient_thresh}
\end{equation}

with $\delta_S = 2$\,\%.
Both conditions are evaluated every $T_{\text{interval}} = 1\,000$\,ms.

\subsubsection{End-to-End Latency Model}

The total latency from sensor event to database record is:

\begin{equation}
	L_{\text{total}} = T_{\text{acq}} + T_{\text{proc}} + T_{\text{serial}} +
	T_{\text{TLS}} + T_{\text{RTT}}
	\label{eq:latency}
\end{equation}

where:
\begin{itemize}
	\item $T_{\text{acq}}$: sensor acquisition time (DHT22: up to 2\,000\,ms;
	      MAX30102: continuous FIFO drain, $\approx$25\,ms per sample at default rate).
	\item $T_{\text{proc}}$: local computation ($<$5\,ms per SpO\textsubscript{2}
	      update; $<$1\,ms for BPM average).
	\item $T_{\text{serial}}$: JSON string build ($<$1\,ms).
	\item $T_{\text{TLS}}$: TLS handshake on first connection (1--3\,s);
	      session reuse on subsequent requests ($<$5\,ms).
	\item $T_{\text{RTT}}$: HTTPS round-trip to Firebase
	      ($\approx$100--300\,ms over 2.4\,GHz Wi-Fi).
\end{itemize}

\subsubsection{Key Assumptions}

\begin{itemize}
	\item 32-bit aligned float reads and writes are atomic on the ESP32
	      architecture, making inter-core sharing of \texttt{float} globals safe for
	      single-writer/single-reader access without a mutex.
	\item Server-side timestamps (\texttt{".sv":"timestamp"}) eliminate the need
	      for an RTC on the device.
	\item A finger is considered present when the IR raw value exceeds 50\,000 counts;
	      below this threshold BPM and SpO\textsubscript{2} are reset to zero.
\end{itemize}

% -------------------------------------------------------
% TODO (Backend Team): add alert threshold model, API
% design rationale, and authentication flow.
% -------------------------------------------------------
\subsection{Backend Subsystem Methodology}
The backend logic is centered on a \textbf{hybrid alerting engine}. Fixed thresholds are applied to vital signs for immediate safety (e.g., HR > 120 or SpO2 < 92\%), while a statistical Z-score model identifies significant deviations from the patient's individualized baseline. The system utilizes \textbf{tRPC} to expose these operations as type-safe procedures. The \textbf{Authentication Flow} uses Firebase ID tokens verified on every request, with \textbf{Role-Based Access Control} (RBAC) determining visibility of clinical data based on staff assignments.

\begin{figure}[H]
    \centering
    \fbox{Figure Placeholder: Flowchart of the Hybrid Alerting Logic showing the parallel processing of fixed thresholds and Z-score anomalies.}
    \caption{Clinical alert and decision logic flow.}
    \label{fig:alert_logic}
\end{figure}

\subsection{Web Dashboard Subsystem Methodology}
The dashboard utilizes a \textbf{reactive component architecture} built with React Hooks. It maintains a real-time connection to the Firebase Realtime Database using a WebSocket-based subscription model. For security and presence tracking, the methodology includes local inference using \textbf{face-api.js} to detect and identify authorized medical personnel. Visualizations are rendered using a rolling buffer of vitals, updated every 1,000\,ms to maintain a high-fidelity trend view.

\subsection{AI Subsystem Methodology}
The AI subsystem formulates the monitoring task as a \textbf{time-series anomaly detection and forecasting problem}.
\begin{itemize}
    \item \textbf{Anomaly Detection:} Calculated via the Z-score $z = \frac{x - \mu}{\sigma}$ over a 20-point sliding window. Values exceeding a magnitude of 2.5 are flagged as clinical outliers.
    \item \textbf{Forecasting:} Implemented via \textbf{Linear Trend Estimation}, where the next 5 intervals are projected using the slope $m$ of the current window.
    \item \textbf{Risk Assessment:} A rule-based classifier synthesizes the anomalies and trends into a discrete clinical status (Stable, Warning, Critical).
\end{itemize}

% ============================================================
\section{Methodology Diagram}
% ============================================================

Figure~\ref{fig:methodology} presents the overall development methodology for the
complete system, from initial design through to deployment.

\begin{figure}[H]
	\centering
	\includegraphics[width=\linewidth]{methodology_diagram.png}
	\caption{Overall development methodology.}
	\label{fig:methodology}
\end{figure}

% flowchart TD
%     %% Main Phases
%     P1["1. System Design"]
%     P2["2. Implementation"]
%     P3["3. Debug Build"]
%     P4["4. Validation"]
%     P5["5. Release Build"]
%     P6["6. Integration"]
%     P7["7. Deployment"]
%
%     %% Sub-descriptions
%     D1["Architecture, math model, pseudocode, schematics"]
%     D2["Firmware, backend, web app, AI model"]
%     D3["--debug build, serial + API logs"]
%     D4["Sensor accuracy, latency, API tests, AI metrics"]
%     D5["Production build, no serial overhead"]
%     D6["End-to-end test:<br>ESP32 -> DB -> Dashboard -> AI"]
%     D7["Flash + configure via AP, go live"]
%
%     %% Main sequential flow
%     P1 --> P2
%     P2 --> P3
%     P3 --> P4
%     P4 --> P5
%     P5 --> P6
%     P6 --> P7
%
%     %% Feedback loop
%     P4 -- "fix & re-implement" --> P2
%
%     %% Dashed connections to sub-descriptions
%     P1 -.-> D1
%     P2 -.-> D2
%     P3 -.-> D3
%     P4 -.-> D4
%     P5 -.-> D5
%     P6 -.-> D6
%     P7 -.-> D7
%
%     %% Styling to visually separate phases from descriptions
%     classDef phase fill:#e0e8f5,stroke:#333,stroke-width:1px,rx:5px,ry:5px;
%     classDef sub fill:#f2f2f2,stroke:#666,stroke-width:1px,rx:3px,ry:3px,font-size:12px;
%
%     class P1,P2,P3,P4,P5,P6,P7 phase;
%     class D1,D2,D3,D4,D5,D6,D7 sub;
